\documentclass[../../main.tex]{subfiles}
\usepackage{amssymb}

\begin{document}

\subsection{Neutrini oltre il Modello Standard}
  \label{subsec:neutrini-oltre-il-modello-standard}
  Come già detto nella sotto sezione
  ~\ref{subsec:il-neutrino-nel-modello-standard}, il Modello Standard
  attribuisce ai neutrini una massa nulla, questa affermazione è però stata
  smentita sia dall'osservazione della loro oscillazione che da studi su
  osservabili cosmologici; il Modello Standard necessita quindi di essere
  esteso.\\
  Il problema della massa dei neutrini si articola in diversi sotto problemi:

  \paragraph{Neutrini di Dirac vs neutrini di Majorana \cite[Capitolo 6]
  {10.1093/acprof:oso/9780198508717.001.0001}}
    L'attuale descrizione dei leptoni implica direttamente l'annullarsi della
    massa dei neutrini tramite il meccanismo di Higgs.
    Le due teorie più accreditate che provano a risolvere questo problema sono
    l'estensione minimale del Modello standard (\textit{minimally
extended Standard Model}) e il Neutrino di Majorana.
    Senza soffermarsi sui dettagli matematici che sono ben oltre gli obbiettivi
    di questa tesi, entrambe le teorie introducono altri tre neutrini (uno per
    generazione) non soggetti alla forza debole (capaci quindi di interagire
    solo gravitazionalmente), in questo modo il meccanismo di Higgs attribuisce
    ai neutrini una massa non nulla.
    La differenza sta in come questi ``nuovi'' neutrini sono legati a quelli già
    noti.
    Esperimenti come CUORE~\cite{ADAMS2022103902} mirano a discriminare queste
    due possibilità.

  \paragraph{Ordinamento degli autostati}
    Gli autostati di massa previsti dalle due teorie appena esposte
    \begin{equation*}
      \bm{n}_L =
      \begin{pmatrix}
        \nu_{1L} \\ \nu_{2L} \\ \nu_{3L}
      \end{pmatrix}
    \end{equation*}
    sono legati a quelli di sapore tramite la matrice
    Pontecorvo-Maki-Nakagawa-Sakata $U^{PMNS}$ dalla relazione
    \begin{equation}
      \bm{\nu}_L = U^{PMNS} \bm{n}_L =
      \begin{pmatrix}
        \Pnue_{L} \\ \Pnum_{L} \\ \Pnut_{L}
      \end{pmatrix}
      \label{eq:massa-sapore}
    \end{equation}
    Gli autostati di massa sono autostati dell'hamiltoniana
    \begin{equation*}
      \mathcal{H} \ket{\nu_k} = E_k \ket{\nu_k}
    \end{equation*}
    ed evolvono quindi come onde piane
    \begin{equation}
      \ket{\nu_k (t)} = e^{-i E_k t} \ket{\nu_k}
      \label{eq:ev-temp-autoket-massa}
    \end{equation}
    Usando due volte l'eq.~\eqref{eq:massa-sapore} nella
    ~\eqref{eq:ev-temp-autoket-massa} possiamo scrivere l'evoluzione temporale
    di un autostato di sapore in funzione degli autostati di sapore stessi
    \begin{equation*}
      \ket{\nu_\alpha} = U^{*}_{\alpha k} e^{-i E_k t} U_{\beta k}
      \ket{\nu_\beta}
    \end{equation*}
    e calcolare la probabilità di transizione tra due sapori come
    \begin{equation*}
      P_{\nu_\alpha \rightarrow \nu_\beta} (t) =
      \left| \braket{\nu_\beta | \nu_\alpha (t)} \right|^2
    \end{equation*}
    Nel limite ultrarelativistico ($v \approx c$), sempre valido visto l'ordine
    di grandezza delle masse dei neutrini,
    $P_{\nu_\alpha \rightarrow \nu_\beta}$ dipende dai termini
    $\Delta m^2_{kj} = m^2_k - m^2_j.$
    Gli esperimenti hanno misurato il segno di $m^2_{12}$, chiamata differenza
    di massa solare, stabilendo che $\nu_1 < \nu_2$, per quanto riguarda
    $m^2_{13}$ e $m^2_{23}$ si sa solo che hanno lo stesso ordine di grandezza;
    la situazione è quindi
    \begin{equation*}
      m^2_{12} \ll \left| m^2_{13} \right| \simeq \left| m^2_{23} \right|
    \end{equation*}
    I due scenari possibili sono quindi la ``gerarchia normale''
    $m_1 < m_2 < m_3$ o la ``gerarchia inversa'' $m_3 < m_1 < m_2.$

  \paragraph{Violazione della simmetria CP}
    La simmetria CP consiste nello scambiare una particella con la usa
    antiparticella: $\Pnu \rightarrow \APnu$; se il sistema non cambia a seguito
    di questa trasformazione allora si dice che rispetta la simmetria CP\@.
    Nel caso dei neutrini questa violazione corrisponderebbe con la condizione
    \begin{equation*}
      P_{\nu_\alpha \rightarrow \nu_\beta} \neq
      P_{\bar{\nu}_\alpha \rightarrow \bar{\nu}_\beta}
    \end{equation*}
    Al momento ci sono diversi esperimenti in corso che provano a rilevare
    questa asimmetria.

\end{document}